\documentclass[11pt]{article}

\usepackage[margin=0.75in]{geometry}
\usepackage{amsmath, amssymb, amsthm}
\usepackage{tcolorbox}

% Theorem Styles
\theoremstyle{definition}
\newtheorem{definition}{Definition}

\theoremstyle{plain}
\newtheorem{theorem}{Theorem}

\theoremstyle{remark}
\newtheorem*{remark}{Integer Arithmetic Rule}

\pagestyle{empty}

\begin{document}

\begin{center}
    {\LARGE \textbf{Math Club Proof Workshop: Direct Proofs \& Divisibility}} \\[4pt]
    {\large Worksheet \& Lab Problem Set}
\end{center}

\vspace{0.5em}
\hrule
\vspace{1em}

\section*{1. Reference Definitions \& Rules}

\begin{definition}[Even]
    An integer \(n\) is \textbf{even} if there exists an integer \(k \in \mathbb{Z}\) such that \(n = 2k\).
\end{definition}

\begin{definition}[Odd]
    An integer \(n\) is \textbf{odd} if there exists an integer \(k \in \mathbb{Z}\) such that \(n = 2k + 1\).
\end{definition}

\begin{definition}[Divisibility]
    For integers \(a\) and \(b\) with \(a \neq 0\), we say \(a\) \textbf{divides} \(b\) (written \(a \mid b\)) if there exists an integer \(k \in \mathbb{Z}\) such that \(b = ak\).
\end{definition}

\begin{remark}
    The sum, difference, and product of any two integers is always an integer (closure under \(+,-,\times\)).
\end{remark}

\vspace{0.5em}
\hrule
\vspace{1em}

\section*{2. Lecture Demo}

\begin{theorem}
    The sum of an even integer and an odd integer is an odd integer.
\end{theorem}

\vspace{0.5em}
\hrule
\vspace{1em}

\section*{3. Proof Lab Problems}
\textit{Note: Use separate scratch paper for non-linear discovery. Once an officer checks your logic, write your final proof in complete English prose below.}

\vspace{1em}

\begin{theorem}[Warmup]
    If \(n\) is an even integer, then \(n^2 + 3n\) is an even integer.
\end{theorem}

\begin{tcolorbox}[colback=white, colframe=black!75, title=Formal Proof (Complete English Sentences)]
\vspace{6.0cm}
\end{tcolorbox}

\newpage

\begin{theorem}[Medium]
    For all integers \(a, b,\) and \(c\), if \(a \mid b\) and \(a \mid c\), then \(a \mid (3b - 2c)\).
\end{theorem}

\begin{tcolorbox}[colback=white, colframe=black!75, title=Formal Proof (Complete English Sentences)]
\vspace{7.0cm}
\end{tcolorbox}

\vspace{1.5em}

\begin{theorem}[Challenge]
    If \(n\) is an odd integer, then \(n^2 - 1\) is divisible by \(8\).
\end{theorem}

\begin{tcolorbox}[colback=white, colframe=black!75, title=Formal Proof (Complete English Sentences)]
\vspace{7.0cm}
\end{tcolorbox}

\end{document}